% ECONOMICS_PROBLEM: {"id": "D91-499", "title": "A unified theory of belief-updating biases", "jel_code": "D91", "source_id": "OP-0067"}
% !TeX program = xelatex
\documentclass[11pt,a4paper]{article}
\usepackage[margin=23mm]{geometry}
\usepackage{fontspec}
\setmainfont{Latin Modern Roman}
\usepackage{amsmath,amssymb,mathtools}
\usepackage{xcolor,enumitem,booktabs,longtable,array,xurl,hyperref}
\definecolor{muted}{HTML}{53636C}
\hypersetup{colorlinks=true,urlcolor=blue,linkcolor=blue}
\setlength{\parindent}{0pt}
\setlength{\parskip}{5pt}
\newcommand{\R}{\mathbb R}
\newcommand{\E}{\mathbb E}
\newcommand{\N}{\mathbb N}
\newcommand{\ind}{\mathbf 1}
\newcommand{\Var}{\operatorname{Var}}
\newcommand{\Cov}{\operatorname{Cov}}
\newcommand{\supp}{\operatorname{supp}}
\DeclareMathOperator*{\argmin}{arg\,min}
\DeclareMathOperator*{\argmax}{arg\,max}
\newcommand{\entrymeta}[3]{{\sffamily\small\color{muted}\textbf{\texttt{#1}}\quad|\quad #2\par #3\par}}
\newcommand{\Problemref}[1]{\texttt{#1}}
\newcommand{\JELref}[2]{\texttt{#2} (JEL \texttt{#1})}
\begin{document}
\section*{A unified theory of belief-updating biases}
\textbf{Problem ID:} \texttt{D91-499}\quad \textbf{JEL:} \texttt{D91}\par
% BEGIN ECONOMICS STATEMENT
\label{problem:OP-0067}
\label{atlas:prob:B7}

\noindent\textbf{Original identifier:} B7 (atlas item 67). \textbf{Classification:} Mechanism discrimination research program. \textbf{Original tractability label:} Hard (editorial, not a complexity result).

\subsection*{Definitions and assumptions}

Let $\theta\in\{0,1\}$ be a state, $p_{i0}\in(0,1)$ a prior, and $\ell_t=\log(f_1(s_t\mid H_t)/f_0(s_t\mid H_t))$ the known log-likelihood ratio of signal $s_t$, where $H_t$ is the presignal information history, and $f_0,f_1$ are positive conditional densities on their common support. The Bayesian benchmark satisfies $L_{it}=L_{i,t-1}+\ell_t$, with $L=\log(p/(1-p))$. A candidate behavioral family may use
\[
L_{it}=a_iL_{i,t-1}+b_i\ell_t+c_i v_t+d_i\operatorname{sgn}(L_{i,t-1})|\ell_t|+\epsilon_{it},
\]
where $v_t$ records prespecified desirability information and $\epsilon_{it}$ has a declared conditional distribution. This is an illustrative descriptive specification, not a definition of all named biases. Elicitation must model risk attitudes, reporting incentives, and selective information acquisition.

\subsection*{Formal research question}

Construct competing structural models in which updating errors arise from signal misperception, attention costs, or preferences over beliefs. For a factorial design $D$ varying signal diagnosticity, information cost, and ego relevance independently, determine whether
\[
P_{D,M,\gamma}=P_{D,M',\gamma'}\Longrightarrow M=M'
\]
holds up to an explicitly defined model equivalence, where $M$ indexes mechanisms and $\gamma$ their parameters. If not, characterize the equivalence classes and an additional intervention that separates them. Compare joint predictions across confirmation, precision, desirability, and prior-sensitivity tasks using a prespecified out-of-sample scoring criterion.

\subsection*{Interpretation and scope}

One mechanism can explain several observed patterns, and several mechanisms can explain the same pattern. Therefore unification should mean a restricted model with cross-task predictions, not a flexible regression assigning a separate coefficient to every bias. Conditional identification requires accurately communicated signal structures; otherwise a measured deviation may reflect reasonable uncertainty about the experiment. Incentivized belief reports identify subjective probabilities only under an appropriate scoring and utility model, or a justified binary-lottery implementation. Overprecision concerns beliefs about uncertainty and cannot generally be reduced to distorted means in a binary-state task. The design must accordingly include tasks measuring distributions or interval coverage. Process measures such as response time are supplementary unless a measurement model relates them to latent computation. A meaningful negative result would demonstrate nonidentification or show that no tested single mechanism achieves the declared prediction tolerance.

\subsection*{Source audit and status}

[S1] formalizes confirmatory misinterpretation; [S2] models motivated self-confidence; [S3] reviews probabilistic reasoning and explicitly omits much motivated-belief research. All are verified, but they do not jointly assert a single unifying conjecture. Existing models already answer particular mechanism questions. The comparative identification exercise above is an editorial research program, with a finite, specified model class; no unrestricted universal theory is promised or certified as open in 2026.

\subsection*{References}

\begin{enumerate}[label={[S\arabic*]},leftmargin=*]

\item Matthew Rabin and Joel L. Schrag (1999). \emph{First Impressions Matter: A Model of Confirmatory Bias}. Quarterly Journal of Economics 114(1), 37--82. \url{https://doi.org/10.1162/003355399555945}\par \textit{Source role:} Model of signal misinterpretation. \textit{Locator:} Abstract and article metadata.

\item Roland Bénabou and Jean Tirole (2002). \emph{Self-Confidence and Personal Motivation}. Quarterly Journal of Economics 117(3), 871--915. \url{https://doi.org/10.1162/003355302760193913}\par \textit{Source role:} Model of motivated belief management. \textit{Locator:} Abstract and article metadata.

\item Daniel J. Benjamin (2019). \emph{Errors in Probabilistic Reasoning and Judgment Biases}. Handbook of Behavioral Economics: Applications and Foundations 1, Volume 2, Chapter 2, 69--186, Elsevier. \url{https://www.sciencedirect.com/science/chapter/handbook/pii/S2352239918300228}\par \textit{Source role:} Review of probabilistic reasoning; motivated beliefs largely outside its scope. \textit{Locator:} Chapter abstract and Section 9 discussion; publisher metadata.

\end{enumerate}

\subsection*{Common conventions: Source mathematical conventions}

Unless an entry states otherwise, agent and firm sets are finite. State and action spaces are measurable, strategies and outcome maps are measurable, probability laws are specified, and expectations exist for the targets under discussion. Infinite-horizon discount factors lie in $(0,1)$ and discounted objectives are finite. Norms, tolerances and complexity input sizes must be specified for computational comparisons. Constants or primitives that appear locally are inputs, not empirical facts asserted by this catalogue.

\textbf{Identification.} A statistical model $m\in\mathcal M$ implies an observable law $P_m$ and target $\theta(m)$. At law $P$, its identified set is
\[
\Theta(P;\mathcal M)=\{\theta(m):m\in\mathcal M,\ P_m=P\}.
\]
Point identification holds when this set is a singleton, or equivalently when $P_{m_1}=P_{m_2}$ implies $\theta(m_1)=\theta(m_2)$ for all compatible models. Sharp bounds mean bounds attained, or approached when the boundary is not attained, by compatible models. Writing this set is a definition, not a solution: the substantive task is to establish informative restrictions and derive or estimate the set. An unrestricted-confounding model may give only trivial bounds.

\textbf{Inference.} Identification, estimability and valid uncertainty quantification are separate questions. Fix $\alpha\in(0,1)$. A confidence procedure $C_n$ over a declared law class $\mathcal P$ has asymptotic uniform coverage at level $1-\alpha$ when
\[
\liminf_{n\to\infty}\inf_{P\in\mathcal P}\Pr_P\{\theta(P)\in C_n\}\geq1-\alpha.
\]
For set-identified targets the coverage event must be defined separately, for example coverage of the full identified set. If an entry uses identification-indexed models instead of uniquely indexed laws, the same coverage convention is applied over those models. Pointwise limits do not establish uniform validity, and asymptotic validity does not guarantee finite-sample precision. Sample sizes, dependence structures and weak-identification sequences are part of each application.

\textbf{Counterfactuals.} $Y_i(a)$ is a potential outcome under a specified intervention, including its timing, implementation and versions. It is not $Y_i$ conditional on voluntarily choosing $a$. A full intervention vector permits interference; shorthand scalar treatment notation does not silently impose no interference. Assignment, selection, attrition, anticipation and measurement must be specified. A claim about an instrument's local effect is not automatically a population policy effect.

\textbf{Economic equilibrium.} A counterfactual model includes best responses, constraints, feasibility, market clearing, state transitions and expectations consistent with the stated information. When several equilibria exist, either a declared selector is an input or the target is set-valued. Comparisons across policy regimes must use the stated selection convention. Reduced-form relationships are not presumed invariant after an equilibrium-changing intervention.

\textbf{Optimization and welfare.} Optimization is over the stated feasible set. If attainment is not established, $\sup$ or $\inf$ and $\varepsilon$-optimal choices replace an asserted maximizer or minimizer. For example, with value $V=\sup_{a\in\mathcal A}W(a)<\infty$, an $\varepsilon$-optimal set is $\{a\in\mathcal A:W(a)\geq V-\varepsilon\}$ for $\varepsilon>0$. Benefits, costs, risk and health or environmental outcomes can be aggregated only using declared units and conversion weights. Interpersonal utility weights, the treatment of fiscal transfers and foreign profits, discounting, and the behavioral welfare criterion are explicit inputs. A comparative-static derivative additionally requires existence and appropriate differentiability of the selected counterfactual.

\textbf{Sources.} Local labels [S1], [S2], and so forth restart in every question. Locators identify the relevant abstract, model, section or result at the level actually checked; a bibliographic or abstract-level verification is not represented as inspection of an entire proof. Working papers are distinguished from later publications and changing versions. Source summaries are paraphrased. All URL checks and editorial formulations refer to the audit date stated above.
% END ECONOMICS STATEMENT
\end{document}
